\chapter{Corrections and proof dependencies}

The following register governs the preserved source papers in the complete reader. It separates false wording, missing hypotheses, and incomplete proof implications. The review was bounded to the current I--VIII manuscripts and selected decisive checks; it does not certify every theorem. The full review notes retain exact source locators.

\section{Endpoint and character conventions}

\textbf{C01 / Volume II, Proposition 13.1: zero depth.}
The number-basis matrix element on $\ell^2(\{2,3,\ldots\})$ omits the $n=1$ term at $k=0$. Its contribution is one under the moment convention $0^0=1$. State the displayed identities for $k\ge1$, or retain the separate zero-eigenvalue atom of the positive contraction. The logarithmic scale generator is used only on $n\ge2$. This is a correction after a stated restriction, not a failure of the positive measure.

\textbf{C02 / Volumes III--V: modulus and conductor.}
For a character modulo $q$, $|\chi(n)|^2=\boldsymbol 1_{(n,q)=1}$. An induced character can have a smaller conductor without changing its actual zero set. The principal character modulo four has conductor one but vanishes at two. Thus the coprimality parameter and squarefree density use the modulus; a conductor formulation requires the character to be primitive modulo that conductor. The consolidated chapters use the modulus convention.

\textbf{C03 / Volume I, Chapter 14: wavelet mean.}
The zero-integral Mertens-wavelet sentence requires $k\ge1$. At $k=0$, its stated integral is one. The subsequent source Mellin theorem already has the needed positive-depth restriction.

\section{Positive spectral models}

\textbf{C04 / Volume III, abstract and claim map: quadrature support.}
Gaussian quadrature nodes lie in the interval spanned by a positive measure; they need not be original discrete spectral points. For $a=\beta=2$, $q=1$, and one quadrature node, the node is $E_1/E_0$. The original spectrum has zero and its first positive point $9/16$. Directly,
\[
0<\frac{E_1}{E_0}<\sum_{n\ge2}n^{-4}
<\frac1{16}+\int_2^\infty x^{-4}\,dx
=\frac5{48}<\frac9{16}.
\]
The node lies in the spectral gap. The actual moment-exact theorem survives this correction. Exact Gaussian nodes here are in $(0,1)$; a tiny floating-point node is not an exact representation of the atom at zero.

\textbf{C05 / Volume III, Corollary 5.2: residue sign.}
The weights in $E_0\lambda_j/(1-z\xi_j)$ are positive. The complex residue at $z=1/\xi_j$ is
\[
\operatorname*{Res}_{z=1/\xi_j}
\frac{E_0\lambda_j}{1-z\xi_j}
=-\frac{E_0\lambda_j}{\xi_j}<0.
\]
The corrected statement is ``positive spectral weights and poles on $(1,\infty)$.''

\textbf{C06 / Volume III, Theorems 6.1--7.1: all moments.}
Weak convergence against bounded continuous functions does not by itself imply convergence of polynomial moments. Chapter 1 supplies the missing implication from Volume II's shifted-energy identity. It also transfers the moments between $U_k=2k n^{-\beta}$ and $u_k=U_k-U_k^2/(4k)$. This closes the proof linkage needed for the fixed-degree Laguerre and determinant arguments.

\textbf{C07 / Volume III, Theorem 11.1: endpoint cusp.}
The original proof informally replaces a discrete measure by a density. Chapter 1 supplies an Abelian proof directly from $E_k\sim A k^{-3/2}$, with no density replacement. The resulting coefficient is the stated $3/(2^{3/2}\pi)$. This repairs the local proof passage conditional on the source energy law and its counting input. It does not provide a new signed arithmetic estimate.

\section{Signed rational models}

\textbf{C08 / Volume IV, Theorem 11.1: denominator degree.}
Hankel nonbreakdown establishes unique normalized rational matching with $\deg P\le m-1$ and $\deg Q\le m$. An ordinary $m$-term simple-packet decomposition additionally follows when $\deg Q=m$ and its $m$ zeros are simple. The source proof does not establish that full-degree implication for every arithmetic packet.

Generic signed moments $(0,1,0,0)$ at $m=2$ have a nonzero Hankel determinant but matching rational function $z$ and denominator one. This illustrates the missing general implication; it is not a counterexample within the M\"obius arithmetic family. Preserve the rational matching result and qualify its simple-packet extension. Positivity, real support, and universal numerical stability are not implied.

\section{Recoupling and deformation}

\textbf{C09 / Volume VII, finite recurrence: sign gauge.}
An irreducible real tridiagonal matrix has nonzero signed off-diagonal entries. Positive entries require a diagonal sign gauge. For the displayed equal-half-shape one-excitation matrix,
\[
R=\begin{pmatrix}1/2&\sqrt3/2\\\sqrt3/2&-1/2\end{pmatrix},
\qquad R\operatorname{diag}(0,2)R^T
=\begin{pmatrix}3/2&-\sqrt3/2\\-\sqrt3/2&1/2\end{pmatrix}.
\]
The displayed matrix therefore cannot satisfy the asserted positive-off-diagonal recurrence in that same gauge. Chapter 3 states the gauge transformation explicitly.

\textbf{C10 / Volumes VII--VIII, pentagon interpretation.}
The product of overlap matrices telescopes for any complete orthonormal bases of one common space. Consistent sign changes, permutations, or a shared incorrect subspace can still give identity. The theorem is valid; the inference that every wrong phase, ordering, or projector must be detected by closure is false. Representation relations, lowest-weight annihilation, pair spectra, and tree eigen-equations must be checked separately. A small closure residual is one consistency check.

\textbf{C11 / Volume VIII, scalar q-Racah family.}
Termination through a denominator parameter $q^{-N}$ alone does not guarantee a positive orthogonal family. Let $N=1$, $q=1/2$, and
\[
\alpha=\beta=1/2,\qquad\gamma=4,\qquad\delta=1.
\]
Then $\gamma q=q^{-N}=2$, all terms needed for $n,y\in\{0,1\}$ have nonzero denominators, yet the two lattice values coincide:
\[
\mu(0)=\mu(1)=3,\qquad R_0=(1,1),\quad R_1=(1,1).
\]
Nondegenerate lattice and admissible positive-weight hypotheses are necessary. Chapter 3 distinguishes a well-defined terminating scalar formula from an orthogonal family.

\textbf{C12 / Volume VIII: scalar-to-matrix comparison.}
The implemented positive Casimir-overlap matrix and the standalone scalar q-Racah fixture are separate constructions. Duality and classical-limit checks of the latter do not establish equality with the former. The exact parameter dictionary, weight normalization, and phase convention remain a comparison obligation. Generic Askey--Wilson eigenfunction identities likewise do not confer positive orthogonality without parameter hypotheses.

\section{Retained open statements and review limits}

The principal M\"obius coherence estimate remains open. Its three ledger conjecture entries are provenance for one recurring target, not three independent conjectures. The higher-excitation q-entropy monotonicity claim is excluded by Volume VIII's own negative control; fixed-kernel entropy nondecrease after measurement remains valid.

Classical recoupling/tridiagonal identifications, strong Schwartz convergence in the heat limit, and full bibliography normalization were not all rederived in this consolidation. Their source dependencies remain explicit. The supplied elementary repairs were self-reviewed; the cusp lemma also received a separate neutral subagent audit. Agreement between reviewers is not a kernel certificate.
